% titel: Blatt 2, A4: Stoffmenge Ethanol aus Volumen und Dichte
% skript: Mortimer Kap. 4.2 (Stoffmenge, Teilchen pro Molekül)

% Teil I: Werkzeuge, kein Lösungsweg.

Eine Stoffmenge ist haeufig nicht über die Masse direkt, sondern über das
Volumen einer Flüssigkeit gegeben. Dichte und Molmasse stellen die Brücke her.
Zusaetzlich geht es um die Stoffmenge einzelner Atome innerhalb der Moleküle.

\begin{werkzeugkasten}[Volumen --- Masse --- Stoffmenge $+$ Atome im Molekül]
\begin{equation}
  m = \varrho \cdot V,\qquad n = \frac{m}{M} = \frac{\varrho V}{M}.
\end{equation}
Einheiten: $\varrho$ in $\mathrm{g/mL}$ und $V$ in $\mathrm{mL}$ liefern $m$ in
$\mathrm{g}$ ($1\ \mathrm{mL} \hat{=}\ 1\ \mathrm{cm^3}$). Dichte für Ethanol:
$\varrho \approx \qty{0.789}{g/mL}$.

Steht ein Molekül wie Ethanol $\mathrm{C_2H_6O}$ für $N_{\mathrm{A}}$
Moleküle, dann ist die Stoffmenge der einzelnen Atomsorten ganzzahlig
hochzuskalieren:
\begin{equation}
  n(\mathrm{C}) = 2\, n(\mathrm{C_2H_6O}),\qquad
  n(\mathrm{H}) = 6\, n(\mathrm{C_2H_6O}),\qquad
  n(\mathrm{O}) = 1\, n(\mathrm{C_2H_6O}).
\end{equation}
Jedes Molekül traegt also $2$ Kohlenstoff-, $6$ Wasserstoff- und $1$
Sauerstoffatom.
\textit{Skript:} Mortimer Kap. 4.2, Übung 2/4.
\end{werkzeugkasten}

\begin{achtung}[Darauf achten]
\begin{itemize}
  \item Dichte statt Molmasse verwenden, um von Volumen auf Masse zu kommen;
        $\mathrm{g}$ nicht mit $\mathrm{mL}$ verwechseln.
  \item Ethanolformel $\mathrm{C_2H_6O}$ (Summenformel) --- die Anzahl der
        Atome je Sorte direkt aus den Indizes.
  \item Für jede Atomsorte getrennt hochskalieren; nicht eine einzige
        ``Gesamt-Atomsumme'' verwechseln mit den Einzelstoffmengen.
  \item Volumen $\mathrm{mL}$ und Dichte $\mathrm{g/mL}$ so kürzen, dass
        Masse in $\mathrm{g}$ übrig bleibt.
\end{itemize}
\end{achtung}
